% !TEX program = lualatex
\documentclass[12pt]{article}
\usepackage{fullpage}
\usepackage{fontspec}
\setmainfont{Linux Libertine O}
% \setmainfont{Times New Roman}
\usepackage{amsmath, amssymb, amsthm, mathtools}
\usepackage{enumitem}
\usepackage{geometry}
\geometry{a4paper, top=2cm, bottom=2cm, left=2.5cm, right=2.5cm}

% Определение команд для распределений
\newcommand{\Normal}{\mathcal{N}}
\newcommand{\Poisson}{\mathrm{Poisson}}

% Окружение для решений
\newenvironment{solution}{\textbf{Решение.}\par}{\hfill\(\square\)}

% Команда для ответа (принимает математическое выражение, без дополнительных \(...\))
\newcommand{\answer}[1]{\textbf{Ответ:} \boxed{#1}}

% Вероятностные обозначения
\newcommand{\Pp}{\mathbb{P}}
\newcommand{\Ee}{\mathbb{E}}
\newcommand{\Var}{\mathrm{Var}}
\newcommand{\Cov}{\mathrm{Cov}}
\newcommand{\ind}{\mathbf{1}}
\newcommand{\dist}{\sim}
\newcommand{\iid}{\stackrel{\mathrm{i.i.d.}}{\sim}}
\newcommand{\disteq}{\stackrel{d}{=}}

% Числовые множества
\newcommand{\N}{\mathbb{N}}
\newcommand{\Z}{\mathbb{Z}}
\newcommand{\Q}{\mathbb{Q}}
\newcommand{\R}{\mathbb{R}}

\title{\textbf{Матстат}\\Домашнее задание №1 к 13.02}
\author{Есиков Сергей Дмитриевич}
\date{\today}

\begin{document}

\maketitle

\begin{enumerate}[label=\textbf{\arabic*.}, left=0pt]

    \item Случайные величины \(U_1, \dots, U_n\) независимы и имеют равномерное распределение
          на \([0, 1]\). Пусть \(U_{(1)}, \dots, U_{(n)}\) их порядковая статистика; обозначим также \(U_{(0)} = 0\) и
          \(U_{(n+1)} = 1\). Докажите, что интервалы между порядковыми статистиками \(V_i := U_{(i)} -
          U_{(i-1)}\) имеют одинаковое распределение при всех \(i = 1, \dots , n+1\) и найдите это распре-
          деление. Верно ли, что \(V_i\) независимые случайные величины?

          \begin{solution}
              \begin{enumerate}
                  \item \(f_{U_i}(x) = \mathbf{1}_{[0, 1]}\)
                  \item \(f_U(X) = f_{U_1, \dots, U_n}(x_1, \dots, x_n) = \prod_1^n f_{U_i}(x) = \mathbf{1}_{[0, 1]^n}\)
                  \item \(\overline{U} = (U_{(1)}, \dots, U_{(n)}, 1)\)
                  \item \(f_{\overline{U}}(X) = f_{\overline{U}}(x_1, \dots, x_n, x_{n+1}) =
                        \begin{cases}
                            n! \cdot \prod_1^n f_{U_i}(x) = n!, &
                            \begin{matrix} 0 \le x_1 < \cdots < x_n \le 1, \\ x_{n+1}=1\end{matrix} \\
                            0,                                  & \text{иначе}
                        \end{cases} \)
                  \item \( A = \left(
                        \begin{matrix}
                                1      & 0      & 0      & \cdots & 0      \\
                                -1     & \ddots & \ddots & \ddots & \vdots \\
                                0      & \ddots & \ddots & \ddots & 0      \\
                                \vdots & \ddots & \ddots & \ddots & 0      \\
                                0      & \cdots & 0      & -1     & 1      \\
                            \end{matrix}\right) \in \mathbb{M}(n+1);
                        \hspace{5pt} \det A = 1; \hspace{5pt}
                        A^{-1} = \left(
                        \begin{matrix}
                                1      & 0      & \cdots & 0      \\
                                1      & \ddots & \ddots & \vdots \\
                                \vdots & \ddots & \ddots & 0      \\
                                1      & \cdots & 1      & 1      \\
                            \end{matrix}\right)\)
                  \item \(V = g(\overline{U}) = A \overline{U}\)
                  \item \(f_V(Y)\) -- такая функция, что \(\int_{g(D)} f_V(Y)dY = \int_{D} f_{\overline{U}}(X)dX\)
                  \item Так как g - диффеоморфизм, можно провести замену переменных \(Y = g(X)\) в интеграле дифформы \\
                        \(\int_{D} f_{\overline{U}}(X)dX = \int_{g(D)} f_{\overline{U}}(g^{-1}(Y)) |\det(J_{g^{-1}})|dY\)
                  \item \(f_V(Y) = f_{\overline{U}}(g^{-1}(Y)) |\det(J_{g^{-1}})|\)
                  \item \(\det(J_{g^{-1}}) = \det(J_{A^{-1}Y}) = \det(A^{-1}) \Rightarrow \det J_{A^{-1}} = 1\)
                  \item \(f_V(Y) = f_{\overline{U}}(A^{-1}(Y)) =
                        \begin{cases}
                            n!, & y_i > 0, \sum_1^{n+1} y_i \\
                            0,  & \text{иначе}
                        \end{cases}\)
                  \item \(\forall i: f_{V_i}(y) = n! \cdot \int_{\substack{v_j > 0 \\ \sum v_j = 1-y}}dv_1 \dots dv_{i-1}dv_{i+1}\dots dv_{n+1} = n! \cdot \frac{(1-y)^{n-1}}{(n-1)!} = n{(1-y)}^{n-1}\) - (объем прямоугольного \((n-1)\) -- симплекса с попправкой на перестановку)
                  \item \(V_i = 1 - \sum_{j \neq i} V_j\) -- явная завивисмость.
              \end{enumerate}
              \answer{\text{Нет, не верно}}
          \end{solution}

    \item В условиях задачи 1 положим \(X_i = f(U_i), i = 1, \dots , n\), где \(f(u) = - \ln(1 - u)\). Каким будет совместное распределение \(X_1, \dots , X_n\)? Проверьте, что порядковую стати-
          стику \(X_{(1)}, \dots , X_{(m)}\) можно найти как \(X_{(i)} = f(U_{(i)})\).

          \begin{solution}
              \begin{enumerate}
                  \item Вместо f буду писать g, чтобы не путаться в плотности и функции.
                  \item По аналогии с предыдущей задачей,
                        \(X = g(U) = \left(\begin{matrix}
                            -\ln(1-U_1) \\
                            \vdots      \\
                            -\ln(1-U_n) \\
                        \end{matrix}\right)\), тогда
                  \item \(f_X = \left(f_U \circ g^{-1}\right) \cdot |\det J_{g^{-1}}|\)
                  \item \(g^{-1}(X) = \left(\begin{matrix}
                                1-\exp(-X_1) \\
                                \vdots       \\
                                1-\exp(-X_n) \\
                            \end{matrix}\right)\)
                  \item \(J_{g^{-1}}(x) = \left(\begin{matrix}
                                \exp(-x_1) & 0      & 0          \\
                                0          & \ddots & 0          \\
                                0          & 0      & \exp(-x_n) \\
                            \end{matrix}\right)\)
                  \item \(\det J_{g^{-1}}(x) = \prod_{1}^{n} \exp(-x_i) = \exp(-\sum_{1}^{n} x_i)\)
                  \item \(f_X(x) = f_U(1-\exp(-x_1), \dots, 1-\exp(-x_n)) \cdot \exp(-\sum_{1}^{n} x_i) = \\ = \exp(-\sum_{1}^{n} x_i) \cdot \mathbf{1}_{g^{-1}(x) \in [0, 1]^n} = \exp(-\sum_{1}^{n} x_i) \cdot \mathbf{1}_{x_i \ge 0, \forall i} = \)
                  \item \(-\ln(1-x)\) - монотонно возрастающая функциия, следовательно \(U_{(1)} \le \cdots \le U_{(n)} \Leftrightarrow g(U_{(1)}) \le \cdots \le g(U_{(n)})\)
              \end{enumerate}
              \answer{f_X(x) = \exp(-\sum_{1}^{n} x_i) \cdot \mathbf{1}_{x_i \ge 0, \forall i}}
          \end{solution}

    \item Пусть \(Z_1, Z_2, \dots\)  н.о.о. случайные величины с экспоненциальным распределением. Положим \(S_0 = 0\), \(S_i = Z_1 + \dots + Z_i\) при \(i \in \N\). Пусть \(V_i\) такие же, как в задаче 1,и при этом независимы от случайной величины \(S_{n+1}\). Покажите, что случайный вектор \((S{n+1}V_1, S_{n+1}V_2, \dots , S_{n+1}V_{n+1})\) имеет такое же распределение, как \((Z_1, Z_2, . . . , Z_{n+1})\).

          \begin{solution}
              \begin{enumerate}
                  \item \(Z = \left(\begin{matrix}
                            Z_1 \\ \vdots \\ Z_{n+1}
                        \end{matrix}\right)\)
                  \item \(f_Z(X) = \mathbf{1}_{{[0, +\infty]}^{n+1}}\cdot\lambda^{n+1} e^{-\lambda \sum_{1}^{n+1}x_i}\)
                  \item \(U \sim Uniform({[0, 1]}^{n+1})\)
                  \item \(f_U = \mathbf{1}_{{[0, 1]^{n+1}}}\)
                  \item \(Z = - \frac{1}{\lambda} \ln{U}\)
                  \item
              \end{enumerate}
          \end{solution}
\end{enumerate}

\end{document}

